\documentclass[12pt,a4paper]{article}
\usepackage[left=1.5cm,right=1.5cm,top=2cm,bottom=2.5cm]{geometry}
\usepackage{multicol}
\usepackage{amsmath,amssymb,amsfonts,latexsym,mathtools}
\usepackage[spanish,es-tabla]{babel}
\usepackage[utf8]{inputenc}
\usepackage[T1]{fontenc}
\usepackage{answers}

% Configuración del paquete answers
\Newassociation{sol}{Solution}{ans}
\newtheorem{ex}{Ejercicio}

\begin{document}

% Encabezado Institucional
\begin{center}
    {\bf\Large Universidad de La Frontera}\\[2pt]
    {\small Departamento de Matemática y Estadística}\\[4pt]
    \hrule height 1.5pt
    \vspace{0.3cm}
    {\small\textbf{Asignatura:} Trigonometría y Geometría Analítica \hfill \textbf{Carrera:} Pedagogía en Matemática}
\end{center}

\vspace{0.2cm}
\begin{center}
    {\Large\bfseries Guía Complementaria N°1: Puntos y Rectas}
\end{center}
\vspace{0.3cm}

\sffamily\small

\begin{itemize} 
\item[{\textbf A)}] \textbf{Resuelva los siguientes problemas relativos a puntos, distancias, división de segmentos y rectas en el plano:}

\Opensolutionfile{ans}[ans1] 

\begin{ex}
Sabiendo que el punto $(9, 2)$ divide al segmento que determinan los puntos $P_1(6, 8)$ y $P_2(x_2, y_2)$ en la relación $r = 3/7$, hallar las coordenadas de $P_2$.
\begin{sol} $(16, -12)$. \end{sol}
\end{ex}

\begin{ex}
Hallar las coordenadas de los vértices de un triángulo sabiendo que las coordenadas de los puntos medios de sus lados son $(-2, 1)$, $(5, 2)$ y $(2, -3)$.
\begin{sol} $(1, 6)$, $(9, -2)$, $(-5, -4)$. \end{sol}
\end{ex}

\begin{ex}
Hallar las coordenadas de los vértices de un triángulo cuyas coordenadas de los puntos medios de sus lados son $(3, 2)$, $(-1, -2)$ y $(5, -4)$.
\begin{sol} $(-3, 4)$, $(9, 0)$, $(1, -8)$. \end{sol}
\end{ex}

\begin{ex}
Demostrar analíticamente que las rectas que unen los puntos medios de los lados adyacentes del cuadrilátero $A(-3, 2)$, $B(5, 4)$, $C(7, -6)$ y $D(-5, -4)$ forman otro cuadrilátero cuyo perímetro es igual a la suma de las diagonales del primero.
\begin{sol} Demostración analítica verificando que los lados del nuevo cuadrilátero son paralelos e iguales a la mitad de las diagonales $AC$ y $BD$. \end{sol}
\end{ex}

\begin{ex}
Demostrar que las rectas que unen los puntos medios de dos lados de los triángulos son paralelas al tercer lado e iguales a su mitad.
\begin{sol} Demostración directa por fórmula de pendiente y distancia entre dos puntos. \end{sol}
\end{ex}

\begin{ex}
Hallar dos puntos de la recta $5x - 12y + 15 = 0$ cuya distancia a $3x + 4y - 12 = 0$ sea $3$.
\begin{sol} $\left(\frac{33}{7}, \frac{45}{14}\right)$, $\left(-\frac{12}{7}, \frac{15}{28}\right)$. \end{sol}
\end{ex}

\begin{ex}
Hallar las ecuaciones de las paralelas a la recta $8x - 15y + 34 = 0$ que distan $3$ unidades del punto $(-2, 3)$.
\begin{sol} $8x - 15y + 112 = 0$, $8x - 15y + 10 = 0$. \end{sol}
\end{ex}

\begin{ex}
Hallar la ecuación de la recta que pasa por el punto $(2, 3)$ y cuya abscisa en el origen es el doble que la ordenada en el origen.
\begin{sol} $x + 2y - 8 = 0$. \end{sol}
\end{ex}

\begin{ex}
Hallar un punto de la recta $3x + y + 4 = 0$ que equidista de los puntos $(-5, 6)$ y $(3, 2)$.
\begin{sol} $(-2, 2)$. \end{sol}
\end{ex}

\begin{ex}
Hallar las ecuaciones de las rectas que pasan por el punto $(1, -6)$ y cuyo producto de coordenadas en el origen es $1$.
\begin{sol} $9x + y - 3 = 0$, $4x + y + 2 = 0$. \end{sol}
\end{ex}

\begin{ex}
Hallar la ecuación de la recta de abscisa en el origen $-3/7$ y que es perpendicular a la recta $3x + 4y - 10 = 0$.
\begin{sol} $28x - 21y + 12 = 0$. \end{sol}
\end{ex}

\begin{ex}
Hallar la ecuación de la perpendicular a la recta $2x + 7y - 3 = 0$ en su punto de intersección con $3x - 2y + 8 = 0$.
\begin{sol} $7x - 2y + 16 = 0$. \end{sol}
\end{ex}

\Closesolutionfile{ans}

\vspace{0.3cm}
\item[{\textbf B)}] \textbf{Ecuaciones de rectas bajo condiciones específicas y familias de rectas:}

\Opensolutionfile{ans}[ans2] 

\begin{ex}
Hallar la ecuación de la recta que pasa por el punto de intersección de las rectas $3x - 5y + 9 = 0$ y $4x + 7y - 28 = 0$ y cumple la condición de pasar por el punto $(-3, -5)$.
\begin{sol} $13x - 8y - 1 = 0$. \end{sol}
\end{ex}

\begin{ex}
Hallar la ecuación de la recta que pasa por el punto de intersección de las rectas $3x - 5y + 9 = 0$ y $4x + 7y - 28 = 0$ y cumple la condición de pasar por el punto $(4, 2)$.
\begin{sol} $38x + 87y - 326 = 0$. \end{sol}
\end{ex}

\begin{ex}
Hallar la ecuación de la recta que pasa por el punto de intersección de las rectas $3x - 5y + 9 = 0$ y $4x + 7y - 28 = 0$ y es paralela a la recta $2x + 3y - 5 = 0$.
\begin{sol} $82x + 123y - 514 = 0$. \end{sol}
\end{ex}

\begin{ex}
Hallar la ecuación de la recta que pasa por el punto de intersección de las rectas $3x - 5y + 9 = 0$ y $4x + 7y - 28 = 0$ y es perpendicular a la recta $4x + 5y - 20 = 0$.
\begin{sol} $205x - 164y + 95 = 0$. \end{sol}
\end{ex}

\begin{ex}
Hallar la ecuación de la recta que pasa por el punto de intersección de las rectas $3x - 5y + 9 = 0$ y $4x + 7y - 28 = 0$ y tiene iguales coordenadas en el origen.
\begin{sol} $41x + 41y - 197 = 0$, $120x - 77y = 0$. \end{sol}
\end{ex}

\begin{ex}
Hallar la ecuación de la recta que pasa por el punto de intersección de las rectas $x - 3y + 1 = 0$ y $2x + 5y - 9 = 0$ y cuya distancia al origen es $2$.
\begin{sol} $x - 2 = 0$, $3x + 4y - 10 = 0$. \end{sol}
\end{ex}

\begin{ex}
Hallar la ecuación de la recta que pasa por el punto de intersección de las rectas $x - 3y + 1 = 0$ y $2x + 5y - 9 = 0$ y cuya distancia al origen es $\sqrt{5}$.
\begin{sol} $2x + y - 5 = 0$. \end{sol}
\end{ex}

\Closesolutionfile{ans}

\vspace{0.3cm}
\item[{\textbf C)}] \textbf{Determine el valor del parámetro $K \in \mathbb{R}$ para que se cumpla cada condición dada:}

\Opensolutionfile{ans}[ans3] 

\begin{ex}
La recta $2x + 3y + K = 0$ forme con los ejes coordenados un triángulo de área $27$ unidades de superficie.
\begin{sol} $K = \pm 18$. \end{sol}
\end{ex}

\begin{ex}
La recta de ecuación $2x + 3Ky - 13 = 0$ pase por el punto $(-2, 4)$.
\begin{sol} $K = 17/12$. \end{sol}
\end{ex}

\begin{ex}
La recta de ecuación $3x - Ky - 8 = 0$ forme un ángulo de $45^\circ$ con la recta $2x + 5y - 17 = 0$.
\begin{sol} $K = 7$, $K = -9/7$. \end{sol}
\end{ex}

\begin{ex}
La recta $(2 + K)x - (3 - K)y + 4K + 14 = 0$ pase por el punto $(2, 3)$.
\begin{sol} $K = -1$. \end{sol}
\end{ex}

\begin{ex}
La recta $Kx + (3 - K)y + 7 = 0$ tenga pendiente igual a $7$.
\begin{sol} $K = 7/2$. \end{sol}
\end{ex}

\begin{ex}
La distancia de la recta $5x - 12y + 3 + K = 0$ al punto $(-3, 2)$ sea, en valor absoluto, igual a $4$.
\begin{sol} $K = -16$, $K = 88$. \end{sol}
\end{ex}

\Closesolutionfile{ans}

\end{itemize}

\newpage
\small
\begin{center}
    {\bfseries\Large Soluciones - Guía Complementaria N°1: Puntos y Rectas}\\[10pt]
\end{center}

\textbf{Sección A}
\begin{multicols}{2}
\input{ans1}
\end{multicols}

\vspace{0.4cm}
\textbf{Sección B}
\begin{multicols}{2}
\input{ans2}
\end{multicols}

\vspace{0.4cm}
\textbf{Sección C}
\begin{multicols}{2}
\input{ans3}
\end{multicols}

\end{document}